\documentclass{article}
\usepackage[margin=1in]{geometry}

\usepackage{hyperref}
\usepackage{url}

\usepackage[utf8]{inputenc}
\usepackage[T1]{fontenc}
\usepackage{booktabs}
\usepackage{amsfonts}
\usepackage{nicefrac}
\usepackage{microtype}
\usepackage{xcolor}

\usepackage{amsmath}
\usepackage{amssymb}
\usepackage{mathtools}
\usepackage{amsthm}

\usepackage[inline]{enumitem}
\setitemize{noitemsep,topsep=0pt,parsep=0pt,partopsep=0pt,leftmargin=*}
\setenumerate{noitemsep,topsep=0pt,parsep=0pt,partopsep=0pt,leftmargin=*}

\usepackage[capitalize,noabbrev]{cleveref}
\crefname{section}{\S}{\S\S}
\crefname{algorithm}{Algorithm}{Algorithms}

\usepackage{latexsym}
\usepackage{mleftright}
\usepackage{xspace}

\usepackage{stmaryrd}
\usepackage{bbm}

\usepackage{tikz}
\usetikzlibrary{arrows,arrows.meta,shapes,automata,positioning}

\usepackage{algorithm}
\usepackage{algorithmic}
\newcommand{\RETURN}{\STATE \textbf{return} }

%%%%%%%%%%%%%%%%%%%%%%%%%%%%%%%%%%%%%%%%%%%%%%%%%%%%%%%%%%
% MACROS (from paper/macros.tex)
%%%%%%%%%%%%%%%%%%%%%%%%%%%%%%%%%%%%%%%%%%%%%%%%%%%%%%%%%%

\colorlet{MacroColor}{black}
\newcommand{\mymacro}[1]{#1}

\newcommand{\defn}[1]{\textbf{#1}}
\newcommand{\defeq}{\mathrel{\overset{\raisebox{-0.25ex}{\textnormal{\tiny def}}}{=}}}

\newcommand{\card}[1]{\left|#1\right|}
\newcommand{\set}[1]{\left\{ #1 \right\}}

\newcommand{\bigo}{{\mymacro{ \mathcal{O}}}}
\newcommand{\kleene}[1]{#1^*}

%%%%%%%%%%%%%%%%%%%%%%%%%%%%%%%%%%%%%%%%%%%%%%%%%%%%%%%%%%
% PAPER-SPECIFIC NOTATION
%%%%%%%%%%%%%%%%%%%%%%%%%%%%%%%%%%%%%%%%%%%%%%%%%%%%%%%%%%

\definecolor{CharacterColor}{RGB}{98,115,19}
\definecolor{TokenColor}{RGB}{140,10,89}
\definecolor{StormBlue}{RGB}{42,90,120}

\newcommand{\preCover}{{\color{TokenColor}\mathcal{P}}}
\newcommand{\Quotient}{{\color{TokenColor}\mathcal{Q}}}
\newcommand{\Remainder}{{\color{TokenColor}\mathcal{R}}}

\newcommand{\basisSet}[1]{\langle #1\rangle}

\newcommand{\tfont}[1]{{\color{TokenColor}\texttt{#1}}}
\newcommand{\charfont}[1]{{\color{CharacterColor}\strut\texttt{#1}}}

\newcommand{\transFnSym}[0]{f}
\newcommand{\transFn}{\ensuremath{{\color{CharacterColor}\transFnSym}}\xspace}

\newcommand{\srcSym}{{\color{TokenColor}x}}
\newcommand{\srcSymp}{{\color{TokenColor}x^\prime}}
\newcommand{\srcStr}{\ensuremath{{\color{TokenColor}\boldsymbol{x}}}\xspace}
\newcommand{\tgtSym}{{\color{CharacterColor}y}}
\newcommand{\tgtStr}{{\color{CharacterColor}\boldsymbol{y}}}

\newcommand{\srcAlphabet}[0]{{\color{TokenColor}\ensuremath{\mathcal{X}}}\xspace}
\newcommand{\tgtAlphabet}[0]{{\color{CharacterColor}\ensuremath{\mathcal{Y}}}\xspace}
\newcommand{\srcStrings}[0]{\ensuremath{\srcAlphabet^*}\xspace}
\newcommand{\tgtStrings}[0]{\ensuremath{\tgtAlphabet^*}\xspace}

\newcommand{\transST}{{\color{CharacterColor}\texttt{f}}}

\newcommand{\States}{\mymacro{\mathsf{S}}}
\newcommand{\Transitions}{\mymacro{\mathsf{T}}}
\newcommand{\InitialStates}[0]{\mymacro{\mathsf{I}}}
\newcommand{\FinalStates}{\mymacro{\mathsf{F}}}

\newcommand{\state}{\mymacro{s}}
\newcommand{\statep}{\mymacro{s^\prime}}

\newcommand{\frontier}{\mathcal{F}}
\newcommand{\powerstate}{S}
\newcommand{\powerstatep}{\powerstate^{\prime}}

\newcommand{\Paths}{\mymacro{\Pi}}

\newcommand{\isCylinder}{\mbox{\texttt{{\color{StormBlue}is\_cylinder}}}}
\newcommand{\run}{\mbox{\texttt{{\color{StormBlue}run}}}}
\newcommand{\step}{\mbox{\texttt{{\color{StormBlue}step}}}}
\newcommand{\isLive}{\mbox{\texttt{{\color{StormBlue}is\_live}}}}

\newcommand{\SafeStates}{\States_{\mathsf{safe}}}
\newcommand{\SafePower}{\mathcal{G}_{\mathsf{safe}}}
\newcommand{\SafePowerK}{\mathcal{G}_{\mathsf{safe}}^{(k)}}

\newcommand{\Succ}{\mathrm{Succ}}

%%%%%%%%%%%%%%%%%%%%%%%%%%%%%%%%%%%%%%%%%%%%%%%%%%%%%%%%%%

\theoremstyle{plain}
\newtheorem{theorem}{Theorem}[section]
\newtheorem{proposition}[theorem]{Proposition}
\newtheorem{lemma}[theorem]{Lemma}
\newtheorem{corollary}[theorem]{Corollary}
\newtheorem{definition}[theorem]{Definition}
\newtheorem{remark}[theorem]{Remark}
\newtheorem{example}[theorem]{Example}

\title{Revising the Finite Decomposition Lemma\\(Lemma~6.1)}
\author{}
\date{}

\begin{document}
\maketitle


%%%%%%%%%%%%%%%%%%%%%%%%%%%%%%%%%%%%%%%%%%%%%%%%%%%%%%%%%%
\section{Current Lemma and Issues}
\label{sec:current}
%%%%%%%%%%%%%%%%%%%%%%%%%%%%%%%%%%%%%%%%%%%%%%%%%%%%%%%%%%

\begin{lemma}[Current Lemma 6.1]
\label{lemma:current}
Let $\transFn\colon \srcStrings\to \tgtStrings$ be a function realized by an
input-determinized transducer $\transST$.
Let $\Paths_{\tgtStr}$ be the set of all paths that emit exactly $\tgtStr$
\begin{equation}
\Paths_{\tgtStr} \defeq  \{ \state_0 \xrightarrow{\srcSym_1:\tgtSym_1}\state_1 \cdots \state_{N-1} \xrightarrow{\srcSym_N:\tgtSym_N}\state_N \mid \state_0\in \InitialStates, \tgtSym_1\cdots\tgtSym_N=\tgtStr\}.
\end{equation}
Let $\States_{\tgtStr}$ be the set of the accepting states of the paths in $\Paths_{\tgtStr}$. For a given target string $\tgtStr\in\tgtStrings$, the precover $\preCover(\tgtStr)$ admits a finite decomposition---i.e., the remainder $\Remainder(\tgtStr)$ and quotient $\Quotient(\tgtStr)$ are both finite sets---if the following two conditions hold:
\begin{enumerate}
\item [(i)] \textbf{Finite pre-image:} The set of paths $\Paths_{\tgtStr}$ is finite -- i.e., there are no $\varepsilon$-cycles on the output on any path that emits $\tgtStr$.
\item [(ii)] \textbf{Safety:} Every state $\state\in\States_{\tgtStr}$ is \defn{safe}, defined inductively, as the smallest set satisfying the following conditions:
\begin{enumerate}
\item [(a)] \textbf{Base case 1 (universality):} $\state$ is universal.
\item [(b)] \textbf{Base case 2 (finite closure):} The language accepted by starting the transducer $\transST$ at $\state$ is finite.
\item [(c)] \textbf{Recursive step:} For all transitions $\state\xrightarrow{\srcSym:\tgtSym}\statep$ the state $\statep$ is safe.
\end{enumerate}
\end{enumerate}
\end{lemma}

\paragraph{Issues.}
(1)~Algorithms compute \emph{frontiers} (sets of reachable states), not
individual paths---$\Paths_{\tgtStr}$ is unnecessary indirection.
(2)~$\States_{\tgtStr}$ are terminal states of paths emitting $\tgtStr$, not
accepting states of the transducer---the name is misleading.
(3)~The algorithm checks \emph{collective} ip-universality of powerstates,
which is strictly weaker than requiring every individual state to be universal.
(4)~The safety condition has no associated algorithm for checking it.


%%%%%%%%%%%%%%%%%%%%%%%%%%%%%%%%%%%%%%%%%%%%%%%%%%%%%%%%%%
\section{Revised Formulation}
\label{sec:revised}
%%%%%%%%%%%%%%%%%%%%%%%%%%%%%%%%%%%%%%%%%%%%%%%%%%%%%%%%%%

\begin{definition}[Frontier and successors]
\label{def:frontier}
The \defn{frontier} $\frontier(\tgtStr) \defeq \{ \state \in \States \mid
\exists\, \state_0 \in \InitialStates,\;
\state_0 \xrightarrow{\srcSym_1:\tgtSym_1} \cdots \xrightarrow{\srcSym_N:\tgtSym_N} \state,\;
\tgtSym_1 \cdots \tgtSym_N = \tgtStr \}$
is the set of states reachable from $\InitialStates$ along arcs whose outputs
concatenate to $\tgtStr$.  For $\powerstate \subseteq \States$ and $\srcSym \in \srcAlphabet$,
$\Succ(\powerstate, \srcSym) \defeq
\{ \statep \mid \exists\, \state \in \powerstate :\;
\state \xrightarrow{\srcSym:\tgtSym} \statep \}$.
\end{definition}


%%%%%%%%%%%%%%%%%%%%%%%%%%%%%%%%%%%%%%%%%%%%%%%%%%%%%%%%%%
\section{Safety}
\label{sec:safety}
%%%%%%%%%%%%%%%%%%%%%%%%%%%%%%%%%%%%%%%%%%%%%%%%%%%%%%%%%%

We give two versions of the safety condition---on individual states and on
powerstates---with the same structure: identify base cases, then check that
every path eventually reaches one.

%----------------------------------------------------------
\subsection{Individual-state safety}
\label{sec:level1}
%----------------------------------------------------------

\begin{definition}[Safe state]
\label{def:safe-state}
$\SafeStates \subseteq \States$ is the smallest set such that $\state \in \SafeStates$ if:
\begin{enumerate}[label=(\alph*)]
\item $L_{\mathrm{inp}}(\state) = \srcStrings$ \quad (ip-universal), or
\item $L(\transST_\state)$ is finite \quad (no cycle reachable from $\state$), or
\item every transition $\state \xrightarrow{\srcSym:\tgtSym} \statep$ has $\statep \in \SafeStates$.
\end{enumerate}
\end{definition}

A state is unsafe iff it can reach a cycle among non-base-case states.  This yields a linear-time algorithm:

\begin{algorithm}[ht]
\caption{$\textsc{ComputeSafeStates}(\transST)$}
\label{alg:safe-states}
\begin{algorithmic}[1]
\STATE $B \leftarrow \{\, \state \mid L_{\mathrm{inp}}(\state) = \srcStrings \,\}
  \;\cup\; \{\, \state \mid \text{no cycle reachable from } \state \,\}$
  \hfill\COMMENT{base cases}
\STATE $G' \leftarrow$ subgraph induced by $\States \setminus B$
\STATE $\mathit{Cyc} \leftarrow$ states in non-trivial SCCs of $G'$
  \hfill\COMMENT{Tarjan}
\STATE $\mathit{Unsafe} \leftarrow \{ \state \in \States \setminus B \mid
  \state \text{ can reach } \mathit{Cyc} \text{ in } G' \}$
  \hfill\COMMENT{reverse BFS}
\RETURN $\SafeStates \leftarrow \States \setminus \mathit{Unsafe}$
\end{algorithmic}
\end{algorithm}

\begin{proposition}
\label{prop:safe-states}
\cref{alg:safe-states} computes exactly $\SafeStates$ in $O(|\States| + |\Transitions|)$ time,
independent of $\tgtStr$.
\end{proposition}
\begin{proof}
If $\state \notin \SafeStates$, it has an unsafe successor; iterating gives an
infinite path in $G'$, hence a cycle, so $\state \in \mathit{Unsafe}$.
Conversely, if $\state$ reaches a non-trivial SCC $C$ in $G'$, states in $C$
cannot satisfy clause~(c), so by induction $\state \notin \SafeStates$.
\end{proof}


%----------------------------------------------------------
\subsection{Powerset safety}
\label{sec:level2}
%----------------------------------------------------------

Individual-state safety misses collective effects: a powerstate
$\{\state_1, \state_2\}$ can be collectively ip-universal even when neither
$\state_1$ nor $\state_2$ is individually universal.
The powerset analysis lifts safety to sets of states, using the same operations
as the paper's algorithms: $\isCylinder$ for collective ip-universality,
$\step_{\tgtStr}$ for successors.

\begin{definition}[Safe powerset]
\label{def:safe-powerset}
$\SafePower \subseteq 2^{\States}$ is the smallest collection such that $\powerstate \in \SafePower$ if:
\begin{enumerate}[label=(\alph*)]
\item $\isCylinder(\powerstate)$
  \quad (collectively ip-universal), or
\item $\bigcup_{\state \in \powerstate} L(\transST_\state)$ is finite
  \quad (finite closure), or
\item for every $\srcSym \in \srcAlphabet$,\;
  $\step_{\tgtStr}(\powerstate, \srcSym) \in \SafePower$ or
  $\step_{\tgtStr}(\powerstate, \srcSym) = \emptyset$.
\end{enumerate}
\end{definition}

\begin{remark}
\label{rem:levels}
$\powerstate \subseteq \SafeStates \implies \powerstate \in \SafePower$,
so the individual analysis is a sound approximation of the powerset analysis.
The converse does not hold.
\end{remark}

The same SCC algorithm applies, now on the \emph{powerset graph} whose
vertices are powerstates and whose edges come from $\step$.  Since this graph
can be exponentially large, we build it lazily with a budget $k$:

\begin{algorithm}[ht]
\caption{$\textsc{ComputeSafePowersets}(\transST, \text{seeds}, k)$}
\label{alg:safe-powersets}
\begin{algorithmic}[1]
\REQUIRE Seeds: powerstates of interest;\quad $k$: budget ($2^{|\States|}$ for full)
\STATE $\mathit{Exp} \leftarrow \emptyset$;\; $Q \leftarrow \text{seeds}$
\WHILE{$Q \ne \emptyset$ and $|\mathit{Exp}| < k$}
  \STATE $\powerstate \leftarrow Q.\mathrm{pop}()$
  \IF{$\powerstate \in \mathit{Exp}$ or $\neg\,\isLive(\powerstate)$}
    \STATE \textbf{continue}
  \ENDIF
  \STATE $\mathit{Exp}.\mathrm{add}(\powerstate)$
  \IF{$\isCylinder(\powerstate)$ or finite closure}
    \STATE mark as base case; \textbf{continue}
  \ENDIF
  \FOR{$\srcSym \in \srcAlphabet$}
    \STATE $\powerstatep \leftarrow \step_{\tgtStr}(\powerstate, \srcSym)$
    \IF{$\powerstatep \notin \mathit{Exp}$}
      \STATE $Q.\mathrm{add}(\powerstatep)$
    \ENDIF
  \ENDFOR
\ENDWHILE
\STATE \textbf{// SCC analysis}
\STATE $B \leftarrow$ base-case powerstates in $\mathit{Exp}$
\STATE $\mathit{Unsafe} \leftarrow \{ \powerstate \in \mathit{Exp} \setminus B \mid
  \powerstate \text{ reaches a non-trivial SCC or unexplored successor} \}$
\RETURN $\SafePowerK \leftarrow \mathit{Exp} \setminus \mathit{Unsafe}$
\end{algorithmic}
\end{algorithm}

Larger $k$ certifies more powerstates; the result is always sound.

\begin{proposition}[Soundness, monotonicity, completeness]
\label{prop:soundness}
For any budget $k$:
(i)~$\SafePowerK \subseteq \SafePower$ \emph{(sound)};
(ii)~$k \le k' \implies \SafePowerK \subseteq \mathcal{G}_{\mathsf{safe}}^{(k')}$
\emph{(monotone)}.
The algorithm always terminates since the powerset graph has at most $2^{|\States|}$
nodes; at completion ($k \ge 2^{|\States|}$), $\SafePowerK = \SafePower$ and
the analysis is \emph{complete}: the decomposition is finite for $\tgtStr$ if
and only if $\frontier(\tgtStr) \in \SafePower$.
\end{proposition}


%%%%%%%%%%%%%%%%%%%%%%%%%%%%%%%%%%%%%%%%%%%%%%%%%%%%%%%%%%
\section{Revised Lemma and Proof}
\label{sec:lemma}
%%%%%%%%%%%%%%%%%%%%%%%%%%%%%%%%%%%%%%%%%%%%%%%%%%%%%%%%%%

\begin{lemma}[Finite Decomposition --- Revised]
\label{lemma:revised}
Let $\transFn\colon \srcStrings \to \tgtStrings$ be a function realized by a
transducer $\transST$.  The precover $\preCover(\tgtStr)$ admits a finite
decomposition if and only if:
\begin{enumerate}
\item [(i)] \textbf{Finite pre-image:} $\frontier(\tgtStr)$ is
  finite (no $\varepsilon$-output cycles on paths emitting $\tgtStr$).
\item [(ii)] \textbf{Safety:} $\frontier(\tgtStr) \in \SafePower$.
\end{enumerate}
\end{lemma}

\begin{corollary}[Individual-state sufficient condition]
$\frontier(\tgtStr) \subseteq \SafeStates \implies$ finite decomposition.
\end{corollary}

\begin{proof}[Proof of \cref{lemma:revised}]
The decomposition BFS explores a tree of powerstates rooted at
$\powerstate_0 = \frontier(\tgtStr)$, with children $\Succ(\powerstate, \srcSym)$.

\emph{Safety propagates}: if $\powerstate \in \SafePower$ is not a base case,
clause~(c) gives $\Succ(\powerstate, \srcSym) \in \SafePower$ for all $\srcSym$.
By (ii), every powerstate in the BFS is safe.

\emph{The tree is finite}: safe powerstates cannot lie on cycles among
non-base-case powerstates (SCC characterization).  Every BFS path terminates at
a base case, an empty powerstate, or a previously-visited powerstate.  Finite
branching + no infinite paths $\Rightarrow$ finite tree (K\H{o}nig).

\emph{Quotient is finite}: each quotient element corresponds to a node reaching
base case~(a) (collectively ip-universal).  Finite tree $\Rightarrow$ finitely many.

\emph{Remainder is finite}: each remainder contribution comes from a
finite-closure powerstate (base case~(b)), contributing $|\bigcup_{\state \in
\powerstate} L(\transST_\state)| < \infty$ strings.  Finitely many such
powerstates $\Rightarrow$ finite remainder.
\end{proof}


%%%%%%%%%%%%%%%%%%%%%%%%%%%%%%%%%%%%%%%%%%%%%%%%%%%%%%%%%%
\section{Connection to the Algorithm}
\label{sec:connection}
%%%%%%%%%%%%%%%%%%%%%%%%%%%%%%%%%%%%%%%%%%%%%%%%%%%%%%%%%%

The paper's $\isCylinder(\srcStr, \tgtStr)$ already performs a BFS over
powerstates from $\run_{\tgtStr}(\srcStr)$ via $\step_{\tgtStr}$, returning
\textsc{True} iff every reachable powerstate is non-empty and
accepting---exactly base case~(a) of powerset safety.
\cref{alg:safe-powersets} extends this with base case~(b) (finite closure) and
SCC analysis to handle cycles, giving a complete characterization.
Termination of the decomposition BFS is guaranteed when
$\frontier(\tgtStr) \in \SafePower$.

In practice:
\begin{enumerate}
\item Precompute $\SafeStates$ via \cref{alg:safe-states} in $O(|\States| + |\Transitions|)$.
  Check $\frontier(\tgtStr) \subseteq \SafeStates$.
\item If that fails, run \cref{alg:safe-powersets} seeded with $\frontier(\tgtStr)$
  and budget $k$.
\end{enumerate}

\paragraph{BPE example.}
Every state of a BPE transducer is ip-universal (total DFA, all accepting), so
$\SafeStates = \States$ and the decomposition is always finite---matching the
empirical observation that BPE quotients carry all mass with empty remainders.


%%%%%%%%%%%%%%%%%%%%%%%%%%%%%%%%%%%%%%%%%%%%%%%%%%%%%%%%%%
\section{Summary}
%%%%%%%%%%%%%%%%%%%%%%%%%%%%%%%%%%%%%%%%%%%%%%%%%%%%%%%%%%

\begin{center}
\renewcommand{\arraystretch}{1.3}
\begin{tabular}{@{}llll@{}}
\toprule
& \textbf{Current} & \textbf{Level 1} & \textbf{Level 2} \\
\midrule
Object & $\Paths_{\tgtStr}$ (paths) & \multicolumn{2}{l}{$\frontier(\tgtStr)$ (frontier)} \\
Safety on & individual states & individual states & powerstates \\
Algorithm & none & \cref{alg:safe-states} & \cref{alg:safe-powersets} \\
Complexity & --- & $O(|\States|{+}|\Transitions|)$ & budgetable ($k$ powerstates) \\
Precomputable & no & yes & seeded per $\tgtStr$ \\
Strength & \multicolumn{2}{l}{sufficient (individual-state)} & iff at $k{=}2^{|\States|}$ \\
\bottomrule
\end{tabular}
\end{center}

\end{document}
